\subsection{Beschreibung des Bändermodells}
% Alt: Aufgabe 1
\Aufgabe{
  \label{Aufg1}
  Beschreiben Sie ausführlich das Bändermodell von Festkörpern und seine Bedeutung für die Unterscheidung zwischen
  Leiter, Halbleiter und Isolator. Gehen Sie dabei insbesondere auf die folgenden Aspekte ein:
  \begin{itemize}
    \item relevante Bänder und deren Besetzung mit Ladungsträgern.
    \item die Bedeutung der Bandlücke und wie diese die elektrischen Eigenschaften von Materialien beeinflusst.
    \item Einfluss von Temperaturänderungen auf die Leitfähigkeit.
  \end{itemize}
}


\Loesung{
  Das Bändermodell erklärt die Energiezustände, die Elektronen in Festkörpern einnehmen können.
  In einem einzelnen Atom sind die Elektronenzustände diskrete Energielevel.
  Siehe: Abschnitt \ref{sec:Bändermodell}
}